\documentclass[12pt]{article}
\usepackage{hyperref}

\usepackage{graphicx}

\usepackage{mathtools}
\DeclarePairedDelimiter\ceil{\lceil}{\rceil}
\DeclarePairedDelimiter\floor{\lfloor}{\rfloor}


\usepackage{amsmath}
\usepackage{amsfonts}
\usepackage{amssymb}

\newcommand{\Mod}[1]{\ (\mathrm{mod}\ #1)}


\usepackage{listings}
\lstset{language=C++,
                basicstyle=\ttfamily,
                keywordstyle=\color{blue}\ttfamily,
                stringstyle=\color{red}\ttfamily,
                commentstyle=\color{green}\ttfamily,
                morecomment=[l][\color{magenta}]{\#}
}

\usepackage{breqn}
\usepackage[]{algorithm2e}
\usepackage{physics}

\newcommand\fnurl[2]{%
  \href{#2}{#1}\footnote{\url{#2}}%
}

\newcommand{\Four}{\mathcal{F}}
\renewcommand{\(}{\left(}
\renewcommand{\)}{\right)}
%% \renewcommand{\{}{\left{}
%% \renewcommand{\}}{\right}}


\title{Bluestein FFT Notes for rocFFT}
\author{The rocFFT Team}



\def\no#1{\nobreak{#1}} % stop equation breaking in dmath

\begin{document}
\maketitle

\section{Introduction}

This document contains notes for computing FFTs via the Bluestein
algorithm for \texttt{rocFFT}.

\section{Bluestein algorithm}

The Fourier trasnform of $\left\{x_n\right\}_{n=0}^{N-1}$ is
\begin{dmath}
  X_k = \sum_{n=0}^{N-1}x_n \zeta_N^{nk}
\end{dmath}
where $\zeta_N = e^{\frac{-2\pi i}{N}}$.

Now,
\begin{dmath}
  kn
  = kn -\frac{1}{2}\(n^2 + k^2\) +\frac{1}{2}\(n^2 + k^2\)
  = \frac{2kn - n^2 - k^2}{2} + \frac{n^2}{2} + \frac{k^2}{2}
  = -\frac{\(k-n\)^2}{2} + \frac{n^2}{2} + \frac{k^2}{2},
\end{dmath}
so
\begin{dmath}
  X_k
  = \sum_{n=0}^{N-1}x_n \zeta_N^{nk}
  = \sum_{n=0}^{N-1}x_n \zeta_N^{-\frac{\(k-n\)^2}{2} + \frac{n^2}{2} + \frac{k^2}{2}}
  = \zeta_N^{\frac{k^2}{2}}
  \sum_{n=0}^{N-1}x_n \zeta_N^{-\frac{\(k-n\)^2}{2} + \frac{n^2}{2}}
  = \zeta_N^{\frac{k^2}{2}}
  \sum_{n=0}^{N-1}x_n \zeta_N^{\frac{n^2}{2}}\zeta_N^{-\frac{\(k-n\)^2}{2}}
  = \zeta_N^{\frac{k^2}{2}}
 \(x_k\zeta_N^{\frac{k^2}{2}}\) * \(\zeta_N^{-\frac{k^2}{2}} \).
\end{dmath}
That is, the Fourier transform of $x_n$ can be expressed as a
convolution of $x_k\zeta_N^{\frac{k^2}{2}}$ and
$\zeta_N^{-\frac{k^2}{2}}$, with the result multiplied by
$\zeta_N^{\frac{k^2}{2}}$.

\section{Convolution Theorem}

The convolution of $\{a_k\}_{k=0}^{N-1}$ and $\{b_k\}_{k=0}^{N-1}$ may
be computed via a Fourier transform using the convolution theorem,
which states that the Fourier transform of a convolution is an
element-wise multiplication (which is referred to as the Hadamard
product).  However, for length-$N$ convolutions, one gets aliases, as
\begin{dmath}
  \mathcal{F}^{-1}_N(A\circ{}B)_k
  = \sum_{\ell=0}^{\ell=N-1}a_{\ell \Mod N} b_{(k-\ell) \Mod N}
\end{dmath}
which is different than the desired linear convolution, 
\begin{dmath}
  (a*b)_k
  = \sum_{\ell=0}^{\ell=k}a_{\ell} b_{k-\ell}.
\end{dmath}
The extra terms are referred to as aliases, and removing them is
dealiasing.  This can be accomplished by extending the input with
zeros from length $N$ to at least length $2N$.  In fact, we can choose
any length greater or equal to $2N$; one typically chooses the length
to be $2^P$, where
\begin{dmath}
  P = \ceil{\log{2N}}.
\end{dmath}


\end{document}
